\chapter{有界变差函数}
\begin{solution}{variation-triangle-norm}
齐次性、三角不等式由变差相应性质得到。若范数为零，则 \(f(a)=0\)，并且
\(|f(x)-f(a)|\le V_a^b(f)=0\)，所以 \(f=0\)。这验证了正定性。
\end{solution}
\begin{solution}{variation-point-spike}
把 \(c\) 放入分割产生上升和下降各 \(|A|\)，其余增量为零；任意分割的和
不超过 \(2|A|\)，所以答案为 \(2|A|\)。若 \(c=a\)，只有一次离开端点的
变化，答案为 \(|A|\)。
\end{solution}
\begin{solution}{variation-product}
写 \(\Delta(fg)=f(t_j)\Delta g+g(t_{j-1})\Delta f\)，取绝对值并求和，
得每个分割的估计。对分割取上确界即得结论；BV 函数有界保证右侧有限。
\end{solution}
\begin{solution}{variation-lipschitz-composition}
每段满足 \(|\phi(f(t_j))-\phi(f(t_{j-1}))|\le L|\Delta f|\)，求和再取确界。
只需在值域上满足 Lipschitz 条件，不要求 \(\phi\) 在整条实线上有定义。
\end{solution}
\begin{solution}{variation-c-one}
由 \(f(t_j)-f(t_{j-1})=\int_{t_{j-1}}^{t_j}f'\) 得 \(V(f)\le\int|f'|\)。
反过来，对任意分割，中值定理给 \(V(f,P)=\sum|f'(\xi_j)|\Delta t_j\)；
沿网格趋零的分割，这些和趋于连续函数 \(|f'|\) 的积分。故确界至少为该积分。
\end{solution}
\begin{solution}{variation-uniform-limit}
固定一个有限分割，逐点收敛足以保证 \(V(f_n,P)\to V(f,P)\)。由
\(V(f_n,P)\le V(f_n)\) 得 \(V(f,P)\le\liminf V(f_n)\)，再对 \(P\) 取确界。
不交换两个上确界或极限，因而无需一致收敛。
\end{solution}
\begin{solution}{variation-complete}
若 \(f_n\) 为该范数的 Cauchy 列，则
\(\|f_n-f_m\|_\infty\le |f_n(a)-f_m(a)|+V(f_n-f_m)\)，故一致收敛到某个
有界函数 \(f\)。固定 \(n\)，用上一题对 \(f_n-f_m\) 取 \(m\to\infty\)，得
\(V(f_n-f)\le\liminf_mV(f_n-f_m)\)。先固定一个充分大的 \(n\)，Cauchy
条件保证右侧有限，于是 \(V(f)\le V(f-f_n)+V(f_n)<\infty\)，即 \(f\in BV\)。
再给定 \(\varepsilon>0\)，取 \(N\) 使 \(m,n\ge N\) 时该范数距离
小于 \(\varepsilon\)。对固定 \(n\ge N\) 让 \(m\to\infty\)，端点值收敛
与变差下半连续性给
\(|f_n(a)-f(a)|+V(f_n-f)\le\varepsilon\)。所以在所给范数下收敛。
\end{solution}
\begin{solution}{variation-representative}
记修改后的函数为 \(g\)，\(g(c)=w\)，且 \(w\) 介于 \(u,v\) 之间。
给定 \(g\) 的任意分割，先加入 \(c\)，变差和不会减小。记 \(c\) 的相邻
分割点为 \(l,r\)，其他分段的贡献之和为 \(A\)。修改后分割和为
\(A+|w-f(l)|+|f(r)-w|\)。

在原函数的分割中去掉 \(c\)，改为加入 \(l<s<c<t<r\)，得
\[
 V_a^b(f)\ge A+|f(s)-f(l)|+|f(t)-f(s)|+|f(r)-f(t)|.
\]
令 \(s\uparrow c\)、\(t\downarrow c\)，右侧趋于
\(A+|u-f(l)|+|v-u|+|f(r)-v|\)。由于
\(|w-u|+|v-w|=|v-u|\)，三角不等式给
\[
 |w-f(l)|+|f(r)-w|
 \le |u-f(l)|+|v-u|+|f(r)-v|.
\]
因此 \(g\) 的每个分割和都不超过 \(V_a^b(f)\)。取上确界便得
\(V_a^b(g)\le V_a^b(f)\)，不需要预先假定局部变差的跳跃公式。
\end{solution}
\begin{solution}{rs-common-jump}
取网格趋零、且 \(c\) 位于某段内部的分割。该段的 \(\alpha\) 增量为 \(1\)，
其他段为零。在该段选择小于 \(c\) 的标记，和为零；选择大于 \(c\) 的标记，
和为一。两种极限不同，故不存在正文意义的积分。
\end{solution}
\begin{solution}{variation-reciprocal}
每段 \(|1/f(t_j)-1/f(t_{j-1})|\le\eta^{-2}|\Delta f|\)，所以
\(V(1/f)\le\eta^{-2}V(f)\)。不需要 \(f\) 保持同号；远离零才是所用条件。
\end{solution}
\begin{solution}{rs-discrete-summable}
部分和的尾项一致范数不超过 \(\sum_{j>N}|a_j|\)，故定义的级数一致收敛。
每个阶跃的变差为 \(1\)，尾项的变差由逐点下半连续性不超过
\(\sum_{j>N}|a_j|\)。因而 \(\alpha\in BV\)，且部分和在变差上趋于 \(\alpha\)。
有限和的 Stieltjes 积分为 \(\sum_{j\le N}a_jg(c_j)\)。连续被积函数的积分
估计给尾部误差至多 \(\|g\|_\infty\sum_{j>N}|a_j|\)，趋零后即得公式。
\end{solution}
